\documentclass[11pt,a4paper]{article}
\usepackage[utf8]{inputenc}
\usepackage[T1]{fontenc}
\usepackage{geometry}
\usepackage{amsmath}
\usepackage{amsfonts}
\usepackage{amssymb}
\usepackage{listings}
\usepackage{xcolor}
\usepackage{graphicx}
\usepackage{hyperref}
\usepackage{tcolorbox}
\usepackage{booktabs}
\usepackage{array}
\usepackage{longtable}
\usepackage{enumitem}
\usepackage{fancyhdr}
\usepackage{titlesec}
\usepackage{algorithm}
\usepackage{algorithmic}
\usepackage{mathtools}
\usepackage{amsthm}

% Page setup
\geometry{margin=2.5cm}
\pagestyle{fancy}
\fancyhf{}
\fancyhead[L]{\textcolor{blue}{RSA Cryptography Educational Guide}}
\fancyhead[R]{\thepage}
\renewcommand{\headrulewidth}{0.4pt}

% Title formatting
\titleformat{\section}{\Large\bfseries\color{blue}}{}{0em}{}
\titleformat{\subsection}{\large\bfseries\color{darkgray}}{}{0em}{}
\titleformat{\subsubsection}{\normalsize\bfseries\color{darkgray}}{}{0em}{}

% Color definitions
\definecolor{codegreen}{rgb}{0,0.6,0}
\definecolor{codegray}{rgb}{0.5,0.5,0.5}
\definecolor{codepurple}{rgb}{0.58,0,0.82}
\definecolor{backcolour}{rgb}{0.95,0.95,0.92}
\definecolor{definitionbox}{rgb}{0.9,0.9,0.9}
\definecolor{examplebox}{rgb}{0.9,0.95,0.9}
\definecolor{warningbox}{rgb}{1,0.95,0.8}
\definecolor{prosbox}{rgb}{0.9,0.95,0.9}
\definecolor{consbox}{rgb}{0.95,0.9,0.9}
\definecolor{mathbox}{rgb}{0.95,0.95,1}

% Code listing setup
\lstset{
    backgroundcolor=\color{backcolour},   
    commentstyle=\color{codegreen},
    keywordstyle=\color{magenta},
    numberstyle=\tiny\color{codegray},
    stringstyle=\color{codepurple},
    basicstyle=\ttfamily\footnotesize,
    breakatwhitespace=false,         
    breaklines=true,                 
    captionpos=b,                    
    keepspaces=true,                 
    numbers=left,                    
    numbersep=5pt,                  
    showspaces=false,                
    showstringspaces=false,
    showtabs=false,                  
    tabsize=2
}

% Custom boxes
\newtcolorbox{definitionbox}{
    colback=definitionbox,
    colframe=red!50!black,
    leftrule=4pt,
    arc=0pt,
    outer arc=0pt
}

\newtcolorbox{examplebox}{
    colback=examplebox,
    colframe=green!50!black,
    arc=0pt,
    outer arc=0pt
}

\newtcolorbox{warningbox}{
    colback=warningbox,
    colframe=orange!50!black,
    leftrule=4pt,
    arc=0pt,
    outer arc=0pt
}

\newtcolorbox{prosbox}{
    colback=prosbox,
    colframe=green!50!black,
    arc=0pt,
    outer arc=0pt
}

\newtcolorbox{consbox}{
    colback=consbox,
    colframe=red!50!black,
    arc=0pt,
    outer arc=0pt
}

\newtcolorbox{mathbox}{
    colback=mathbox,
    colframe=blue!50!black,
    leftrule=4pt,
    arc=0pt,
    outer arc=0pt
}

% Theorem environments
\newtheorem{theorem}{Theorem}[section]
\newtheorem{lemma}[theorem]{Lemma}
\newtheorem{corollary}[theorem]{Corollary}
\newtheorem{definition}[theorem]{Definition}

% Title and author
\title{\Huge\textbf{RSA Cryptography: A Comprehensive Guide}\\[0.5cm]
\large From Mathematical Theory to C++17 Implementation}
\author{CryptoGL Educational Series}
\date{\today}

\begin{document}

\maketitle

\begin{abstract}
This document provides a comprehensive introduction to RSA (Rivest-Shamir-Adleman) cryptography, one of the most widely used asymmetric encryption algorithms. We begin with the mathematical foundations, including number theory concepts essential for understanding RSA. We then present step-by-step practical examples of RSA encryption and decryption, followed by efficient C++17 implementation strategies. This guide is designed for beginners with basic mathematical knowledge who want to understand both the theory and practical implementation of RSA cryptography.
\end{abstract}

\tableofcontents
\newpage

\section{Introduction to RSA Cryptography}

RSA is an asymmetric cryptographic algorithm named after its creators: Ron Rivest, Adi Shamir, and Leonard Adleman, who first publicly described it in 1977. It is based on the mathematical difficulty of factoring the product of two large prime numbers.

\begin{definitionbox}
\textbf{RSA Cryptography:} An asymmetric cryptographic system that uses a pair of mathematically related keys - a public key for encryption and a private key for decryption. The security of RSA relies on the computational difficulty of factoring large composite numbers.
\end{definitionbox}

\subsection{Key Concepts}

\begin{itemize}
    \item \textbf{Asymmetric Encryption:} Uses different keys for encryption and decryption
    \item \textbf{Public Key:} Can be freely shared and used by anyone to encrypt messages
    \item \textbf{Private Key:} Must be kept secret and is used for decryption
    \item \textbf{Digital Signatures:} Can also be used to create digital signatures
\end{itemize}

\section{Mathematical Foundations}

\subsection{Number Theory Prerequisites}

\subsubsection{Prime Numbers}

\begin{definition}
A \textbf{prime number} is a natural number greater than 1 that has no positive divisors other than 1 and itself.
\end{definition}

\begin{examplebox}
\textbf{Examples of Prime Numbers:}
\begin{itemize}
    \item 2, 3, 5, 7, 11, 13, 17, 19, 23, 29, 31, 37, 41, 43, 47, 53, 59, 61, 67, 71, 73, 79, 83, 89, 97
    \item For RSA, we typically use primes with hundreds or thousands of digits
\end{itemize}
\end{examplebox}

\subsubsection{Greatest Common Divisor (GCD)}

\begin{definition}
The \textbf{greatest common divisor} of two integers $a$ and $b$, denoted $\gcd(a,b)$, is the largest positive integer that divides both $a$ and $b$.
\end{definition}

\begin{mathbox}
\textbf{Euclidean Algorithm for GCD:}
\begin{algorithmic}
\STATE \textbf{Input:} Two positive integers $a$ and $b$
\STATE \textbf{Output:} $\gcd(a,b)$
\WHILE{$b \neq 0$}
    \STATE $r \leftarrow a \bmod b$
    \STATE $a \leftarrow b$
    \STATE $b \leftarrow r$
\ENDWHILE
\RETURN $a$
\end{algorithmic}
\end{mathbox}

\begin{examplebox}
\textbf{Example:} Find $\gcd(48, 18)$
\begin{align*}
48 &= 2 \times 18 + 12 \\
18 &= 1 \times 12 + 6 \\
12 &= 2 \times 6 + 0 \\
\end{align*}
Therefore, $\gcd(48, 18) = 6$
\end{examplebox}

\subsubsection{Modular Arithmetic}

\begin{definition}
For integers $a$, $b$, and $n > 0$, we say $a \equiv b \pmod{n}$ if $n$ divides $(a - b)$. This means $a$ and $b$ have the same remainder when divided by $n$.
\end{definition}

\begin{mathbox}
\textbf{Properties of Modular Arithmetic:}
\begin{align*}
(a + b) \bmod n &= ((a \bmod n) + (b \bmod n)) \bmod n \\
(a \times b) \bmod n &= ((a \bmod n) \times (b \bmod n)) \bmod n \\
(a^b) \bmod n &= ((a \bmod n)^b) \bmod n
\end{align*}
\end{mathbox}

\subsubsection{Euler's Totient Function}

\begin{definition}
Euler's totient function $\phi(n)$ counts the number of integers between 1 and $n$ that are coprime to $n$ (i.e., $\gcd(k,n) = 1$ for $1 \leq k \leq n$).
\end{definition}

\begin{theorem}
If $n = pq$ where $p$ and $q$ are distinct prime numbers, then:
$$\phi(n) = \phi(pq) = (p-1)(q-1)$$
\end{theorem}

\begin{examplebox}
\textbf{Example:} Let $p = 3$ and $q = 5$, then $n = 15$
\begin{align*}
\phi(15) &= \phi(3 \times 5) = (3-1)(5-1) = 2 \times 4 = 8
\end{align*}
The numbers coprime to 15 are: 1, 2, 4, 7, 8, 11, 13, 14 (8 numbers)
\end{examplebox}

\subsubsection{Fermat's Little Theorem}

\begin{theorem}[Fermat's Little Theorem]
If $p$ is a prime number and $a$ is an integer not divisible by $p$, then:
$$a^{p-1} \equiv 1 \pmod{p}$$
\end{theorem}

\subsubsection{Euler's Theorem}

\begin{theorem}[Euler's Theorem]
If $\gcd(a,n) = 1$, then:
$$a^{\phi(n)} \equiv 1 \pmod{n}$$
\end{theorem}

\subsection{Extended Euclidean Algorithm}

The extended Euclidean algorithm not only finds $\gcd(a,b)$ but also finds integers $x$ and $y$ such that:
$$ax + by = \gcd(a,b)$$

\begin{mathbox}
\textbf{Extended Euclidean Algorithm:}
\begin{algorithmic}
\STATE \textbf{Input:} Two positive integers $a$ and $b$
\STATE \textbf{Output:} $\gcd(a,b)$ and integers $x, y$ such that $ax + by = \gcd(a,b)$
\STATE $x_0, x_1 \leftarrow 1, 0$
\STATE $y_0, y_1 \leftarrow 0, 1$
\WHILE{$b \neq 0$}
    \STATE $q \leftarrow \lfloor a/b \rfloor$
    \STATE $r \leftarrow a \bmod b$
    \STATE $x_2 \leftarrow x_0 - q \times x_1$
    \STATE $y_2 \leftarrow y_0 - q \times y_1$
    \STATE $a, b \leftarrow b, r$
    \STATE $x_0, x_1 \leftarrow x_1, x_2$
    \STATE $y_0, y_1 \leftarrow y_1, y_2$
\ENDWHILE
\RETURN $(a, x_0, y_0)$
\end{algorithmic}
\end{mathbox}

\section{RSA Algorithm}

\subsection{Key Generation}

The RSA key generation process involves the following steps:

\begin{enumerate}
    \item Choose two distinct prime numbers $p$ and $q$
    \item Compute $n = pq$ (the modulus)
    \item Compute $\phi(n) = (p-1)(q-1)$
    \item Choose an integer $e$ such that $1 < e < \phi(n)$ and $\gcd(e, \phi(n)) = 1$
    \item Compute $d$ such that $ed \equiv 1 \pmod{\phi(n)}$ (using extended Euclidean algorithm)
\end{enumerate}

\begin{mathbox}
\textbf{RSA Key Components:}
\begin{align*}
\text{Public Key} &= (n, e) \\
\text{Private Key} &= (n, d) \\
\text{where } ed &\equiv 1 \pmod{\phi(n)}
\end{align*}
\end{mathbox}

\subsection{Encryption and Decryption}

\begin{definitionbox}
\textbf{RSA Encryption:} For a message $M$ (where $0 \leq M < n$):
$$C = M^e \bmod n$$

\textbf{RSA Decryption:} For a ciphertext $C$:
$$M = C^d \bmod n$$
\end{definitionbox}

\subsection{Mathematical Proof of Correctness}

\begin{theorem}
If $C = M^e \bmod n$ and $M = C^d \bmod n$, then the RSA algorithm correctly encrypts and decrypts messages.
\end{theorem}

\begin{proof}
We need to prove that $(M^e)^d \equiv M \pmod{n}$.

Since $ed \equiv 1 \pmod{\phi(n)}$, we have $ed = 1 + k\phi(n)$ for some integer $k$.

Therefore:
$$(M^e)^d = M^{ed} = M^{1 + k\phi(n)} = M \times M^{k\phi(n)}$$

By Euler's theorem, if $\gcd(M,n) = 1$, then $M^{\phi(n)} \equiv 1 \pmod{n}$, so:
$$M^{k\phi(n)} \equiv 1^k \equiv 1 \pmod{n}$$

Thus:
$$(M^e)^d \equiv M \times 1 \equiv M \pmod{n}$$

For the case where $\gcd(M,n) \neq 1$, the proof is more complex but still holds.
\end{proof}

\section{Step-by-Step Practical Examples}

\subsection{Example 1: Small Numbers for Understanding}

Let's work through a simple example with small numbers to understand the process.

\begin{examplebox}
\textbf{Step 1: Key Generation}
\begin{itemize}
    \item Choose $p = 3$ and $q = 11$
    \item Compute $n = p \times q = 3 \times 11 = 33$
    \item Compute $\phi(n) = (p-1)(q-1) = 2 \times 10 = 20$
    \item Choose $e = 3$ (since $\gcd(3, 20) = 1$)
    \item Find $d$ such that $3d \equiv 1 \pmod{20}$
\end{itemize}

\textbf{Step 2: Finding the Private Key}
Using extended Euclidean algorithm to find $d$:
\begin{align*}
20 &= 6 \times 3 + 2 \\
3 &= 1 \times 2 + 1 \\
2 &= 2 \times 1 + 0
\end{align*}
Working backwards:
\begin{align*}
1 &= 3 - 1 \times 2 \\
&= 3 - 1 \times (20 - 6 \times 3) \\
&= 7 \times 3 - 1 \times 20
\end{align*}
Therefore, $d = 7$ (since $3 \times 7 = 21 \equiv 1 \pmod{20}$)

\textbf{Keys:}
\begin{align*}
\text{Public Key} &= (n=33, e=3) \\
\text{Private Key} &= (n=33, d=7)
\end{align*}

\textbf{Step 3: Encryption}
Let's encrypt the message $M = 5$:
$$C = 5^3 \bmod 33 = 125 \bmod 33 = 26$$

\textbf{Step 4: Decryption}
Let's decrypt the ciphertext $C = 26$:
$$M = 26^7 \bmod 33 = 8031810176 \bmod 33 = 5$$

The original message is recovered!
\end{examplebox}

\subsection{Example 2: Larger Numbers}

\begin{examplebox}
\textbf{Step 1: Key Generation}
\begin{itemize}
    \item Choose $p = 61$ and $q = 53$
    \item Compute $n = p \times q = 61 \times 53 = 3233$
    \item Compute $\phi(n) = (p-1)(q-1) = 60 \times 52 = 3120$
    \item Choose $e = 17$ (since $\gcd(17, 3120) = 1$)
    \item Find $d$ such that $17d \equiv 1 \pmod{3120}$
\end{itemize}

\textbf{Step 2: Finding the Private Key}
Using extended Euclidean algorithm:
\begin{align*}
3120 &= 183 \times 17 + 9 \\
17 &= 1 \times 9 + 8 \\
9 &= 1 \times 8 + 1 \\
8 &= 8 \times 1 + 0
\end{align*}
Working backwards:
\begin{align*}
1 &= 9 - 1 \times 8 \\
&= 9 - 1 \times (17 - 1 \times 9) \\
&= 2 \times 9 - 1 \times 17 \\
&= 2 \times (3120 - 183 \times 17) - 1 \times 17 \\
&= 2 \times 3120 - 367 \times 17
\end{align*}
Therefore, $d = -367 \equiv 2753 \pmod{3120}$

\textbf{Keys:}
\begin{align*}
\text{Public Key} &= (n=3233, e=17) \\
\text{Private Key} &= (n=3233, d=2753)
\end{align*}

\textbf{Step 3: Encryption}
Let's encrypt the message $M = 123$:
$$C = 123^{17} \bmod 3233 = 855$$

\textbf{Step 4: Decryption}
Let's decrypt the ciphertext $C = 855$:
$$M = 855^{2753} \bmod 3233 = 123$$

The original message is recovered!
\end{examplebox}

\section{Efficient Implementation in C++17}

\subsection{Mathematical Optimizations}

\subsubsection{Fast Modular Exponentiation}

The naive approach of computing $a^b \bmod n$ by first computing $a^b$ and then taking the modulus is computationally infeasible for large numbers. We use the square-and-multiply algorithm (also known as binary exponentiation).

\begin{mathbox}
\textbf{Square-and-Multiply Algorithm:}
\begin{algorithmic}
\STATE \textbf{Input:} $a, b, n$
\STATE \textbf{Output:} $a^b \bmod n$
\STATE $result \leftarrow 1$
\STATE $base \leftarrow a \bmod n$
\WHILE{$b > 0$}
    \IF{$b$ is odd}
        \STATE $result \leftarrow (result \times base) \bmod n$
    \ENDIF
    \STATE $base \leftarrow (base \times base) \bmod n$
    \STATE $b \leftarrow b \div 2$
\ENDWHILE
\RETURN $result$
\end{algorithmic}
\end{mathbox}

\subsubsection{Chinese Remainder Theorem (CRT)}

For decryption, we can use the Chinese Remainder Theorem to speed up computation:

\begin{theorem}[Chinese Remainter Theorem]
If $n = pq$ where $p$ and $q$ are coprime, and we know $M \bmod p$ and $M \bmod q$, then we can efficiently compute $M \bmod n$.
\end{theorem}

For RSA decryption with CRT:
\begin{align*}
M_p &= C^{d_p} \bmod p \text{ where } d_p = d \bmod (p-1) \\
M_q &= C^{d_q} \bmod q \text{ where } d_q = d \bmod (q-1) \\
M &= M_p + p \times ((M_q - M_p) \times p^{-1} \bmod q)
\end{align*}

\subsection{C++17 Implementation}

\subsubsection{Basic RSA Class Structure}

\begin{lstlisting}[language=C++, caption=Basic RSA Class Structure]
#include <iostream>
#include <vector>
#include <random>
#include <chrono>
#include <cstdint>

class RSA {
private:
    // Key components
    uint64_t n;  // modulus
    uint64_t e;  // public exponent
    uint64_t d;  // private exponent
    uint64_t p;  // first prime
    uint64_t q;  // second prime
    
    // Helper functions
    uint64_t gcd(uint64_t a, uint64_t b);
    uint64_t mod_inverse(uint64_t a, uint64_t m);
    uint64_t mod_pow(uint64_t base, uint64_t exp, uint64_t mod);
    bool is_prime(uint64_t n);
    uint64_t generate_prime(uint64_t min, uint64_t max);

public:
    RSA(uint64_t bit_length = 1024);
    uint64_t encrypt(uint64_t message);
    uint64_t decrypt(uint64_t ciphertext);
    std::pair<uint64_t, uint64_t> get_public_key() const;
};
\end{lstlisting}

\subsubsection{Mathematical Helper Functions}

\begin{lstlisting}[language=C++, caption=Mathematical Helper Functions]
// Greatest Common Divisor using Euclidean algorithm
uint64_t RSA::gcd(uint64_t a, uint64_t b) {
    while (b != 0) {
        uint64_t temp = b;
        b = a % b;
        a = temp;
    }
    return a;
}

// Extended Euclidean Algorithm for modular inverse
uint64_t RSA::mod_inverse(uint64_t a, uint64_t m) {
    int64_t m0 = m, t, q;
    int64_t x0 = 0, x1 = 1;
    
    if (m == 1) return 0;
    
    while (a > 1) {
        q = a / m;
        t = m;
        m = a % m;
        a = t;
        t = x0;
        x0 = x1 - q * x0;
        x1 = t;
    }
    
    if (x1 < 0) x1 += m0;
    return x1;
}

// Fast modular exponentiation using square-and-multiply
uint64_t RSA::mod_pow(uint64_t base, uint64_t exp, uint64_t mod) {
    uint64_t result = 1;
    base = base % mod;
    
    while (exp > 0) {
        if (exp & 1) {
            result = (result * base) % mod;
        }
        base = (base * base) % mod;
        exp >>= 1;
    }
    return result;
}
\end{lstlisting}

\subsubsection{Prime Number Generation}

\begin{lstlisting}[language=C++, caption=Prime Number Generation]
// Miller-Rabin primality test
bool RSA::is_prime(uint64_t n) {
    if (n <= 1) return false;
    if (n <= 3) return true;
    if (n % 2 == 0) return false;
    
    // Write n as 2^r * d + 1
    uint64_t d = n - 1;
    uint64_t r = 0;
    while (d % 2 == 0) {
        d /= 2;
        r++;
    }
    
    // Test with first few prime bases
    std::vector<uint64_t> bases = {2, 3, 5, 7, 11, 13, 17, 19, 23, 29, 31, 37};
    
    for (uint64_t base : bases) {
        if (base >= n) continue;
        
        uint64_t x = mod_pow(base, d, n);
        if (x == 1 || x == n - 1) continue;
        
        bool is_composite = true;
        for (uint64_t i = 1; i < r; i++) {
            x = (x * x) % n;
            if (x == n - 1) {
                is_composite = false;
                break;
            }
        }
        if (is_composite) return false;
    }
    return true;
}

// Generate a random prime number
uint64_t RSA::generate_prime(uint64_t min, uint64_t max) {
    std::random_device rd;
    std::mt19937_64 gen(rd());
    std::uniform_int_distribution<uint64_t> dis(min, max);
    
    uint64_t candidate;
    do {
        candidate = dis(gen);
        // Ensure odd number
        if (candidate % 2 == 0) candidate++;
    } while (!is_prime(candidate));
    
    return candidate;
}
\end{lstlisting}

\subsubsection{RSA Key Generation}

\begin{lstlisting}[language=C++, caption=RSA Key Generation]
// Constructor with key generation
RSA::RSA(uint64_t bit_length) {
    uint64_t min_prime = 1ULL << (bit_length / 2 - 1);
    uint64_t max_prime = 1ULL << (bit_length / 2);
    
    // Generate two large prime numbers
    p = generate_prime(min_prime, max_prime);
    do {
        q = generate_prime(min_prime, max_prime);
    } while (q == p);
    
    // Calculate modulus
    n = p * q;
    
    // Calculate Euler's totient function
    uint64_t phi_n = (p - 1) * (q - 1);
    
    // Choose public exponent (commonly 65537)
    e = 65537;
    while (gcd(e, phi_n) != 1) {
        e += 2;
    }
    
    // Calculate private exponent
    d = mod_inverse(e, phi_n);
}
\end{lstlisting}

\subsubsection{Encryption and Decryption}

\begin{lstlisting}[language=C++, caption=Encryption and Decryption]
// Encrypt a message
uint64_t RSA::encrypt(uint64_t message) {
    if (message >= n) {
        throw std::invalid_argument("Message too large for current key size");
    }
    return mod_pow(message, e, n);
}

// Decrypt a ciphertext
uint64_t RSA::decrypt(uint64_t ciphertext) {
    return mod_pow(ciphertext, d, n);
}

// Get public key
std::pair<uint64_t, uint64_t> RSA::get_public_key() const {
    return {n, e};
}
\end{lstlisting}

\subsubsection{Complete Example Usage}

\begin{lstlisting}[language=C++, caption=Complete RSA Example]
#include <iostream>
#include <string>

int main() {
    try {
        // Create RSA instance with 1024-bit keys
        RSA rsa(1024);
        
        // Get public key
        auto [n, e] = rsa.get_public_key();
        std::cout << "Public Key (n, e): (" << n << ", " << e << ")\n";
        
        // Example message
        uint64_t message = 12345;
        std::cout << "Original message: " << message << "\n";
        
        // Encrypt
        uint64_t encrypted = rsa.encrypt(message);
        std::cout << "Encrypted: " << encrypted << "\n";
        
        // Decrypt
        uint64_t decrypted = rsa.decrypt(encrypted);
        std::cout << "Decrypted: " << decrypted << "\n";
        
        // Verify
        if (message == decrypted) {
            std::cout << "RSA encryption/decryption successful!\n";
        } else {
            std::cout << "Error in RSA implementation!\n";
        }
        
    } catch (const std::exception& e) {
        std::cerr << "Error: " << e.what() << "\n";
    }
    
    return 0;
}
\end{lstlisting}

\subsection{Advanced Optimizations}

\subsubsection{Chinese Remainder Theorem Implementation}

\begin{lstlisting}[language=C++, caption=CRT-Optimized Decryption]
// CRT-optimized decryption
uint64_t RSA::decrypt_crt(uint64_t ciphertext) {
    // Precompute CRT parameters
    uint64_t d_p = d % (p - 1);
    uint64_t d_q = d % (q - 1);
    uint64_t q_inv = mod_inverse(q, p);
    
    // Compute using CRT
    uint64_t m_p = mod_pow(ciphertext, d_p, p);
    uint64_t m_q = mod_pow(ciphertext, d_q, q);
    
    // Combine using CRT
    uint64_t h = (q_inv * (m_p - m_q)) % p;
    uint64_t m = m_q + h * q;
    
    return m;
}
\end{lstlisting}

\subsubsection{Big Integer Support}

For production use, you'll need arbitrary-precision arithmetic. Here's how to integrate with a big integer library:

\begin{lstlisting}[language=C++, caption=Big Integer RSA Implementation]
#include "big_integers/BigIntegerLibrary.hh"

class BigIntRSA {
private:
    BigInteger n, e, d, p, q;
    
    BigInteger mod_pow(const BigInteger& base, const BigInteger& exp, 
                      const BigInteger& mod) {
        BigInteger result = 1;
        BigInteger base_mod = base % mod;
        BigInteger exp_copy = exp;
        
        while (exp_copy > 0) {
            if (exp_copy % 2 == 1) {
                result = (result * base_mod) % mod;
            }
            base_mod = (base_mod * base_mod) % mod;
            exp_copy /= 2;
        }
        return result;
    }
    
public:
    BigIntRSA(int bit_length = 2048) {
        // Generate large primes using a proper big integer library
        // This is a simplified version - in practice, use established libraries
        // like GMP, OpenSSL, or Crypto++ for secure prime generation
    }
    
    BigInteger encrypt(const BigInteger& message) {
        return mod_pow(message, e, n);
    }
    
    BigInteger decrypt(const BigInteger& ciphertext) {
        return mod_pow(ciphertext, d, n);
    }
};
\end{lstlisting}

\section{Security Considerations}

\subsection{Key Size Requirements}

\begin{warningbox}
\textbf{Important:} The security of RSA depends on the difficulty of factoring the modulus $n$. As computational power increases, larger key sizes are required.
\end{warningbox}

\begin{table}[h]
\centering
\begin{tabular}{|c|c|c|}
\hline
\textbf{Key Size (bits)} & \textbf{Security Level} & \textbf{Recommended Until} \\
\hline
1024 & 80 bits & 2010 \\
2048 & 112 bits & 2030 \\
3072 & 128 bits & 2040 \\
4096 & 192 bits & 2050+ \\
\hline
\end{tabular}
\caption{RSA Key Size Recommendations}
\end{table}

\subsection{Common Attacks and Mitigations}

\subsubsection{Factorization Attacks}

\begin{itemize}
    \item \textbf{General Number Field Sieve (GNFS):} Most efficient known algorithm for factoring large integers
    \item \textbf{Quadratic Sieve:} Effective for numbers up to about 100 digits
    \item \textbf{Mitigation:} Use sufficiently large key sizes (2048+ bits)
\end{itemize}

\subsubsection{Timing Attacks}

\begin{itemize}
    \item \textbf{Attack:} Measure time taken for encryption/decryption to infer information about the private key
    \item \textbf{Mitigation:} Use constant-time implementations and blinding techniques
\end{itemize}

\subsubsection{Chosen Ciphertext Attacks}

\begin{itemize}
    \item \textbf{Attack:} Exploit properties of RSA to decrypt messages without knowing the private key
    \item \textbf{Mitigation:} Use proper padding schemes (PKCS\#1, OAEP)
\end{itemize}

\subsection{Padding Schemes}

\begin{definitionbox}
\textbf{Padding:} Adding random data to messages before encryption to prevent certain attacks and ensure messages are the correct length.
\end{definitionbox}

\subsubsection{PKCS\#1 v1.5 Padding}

\begin{lstlisting}[language=C++, caption=PKCS#1 v1.5 Padding]
std::vector<uint8_t> pkcs1_pad(const std::vector<uint8_t>& message, 
                               size_t block_size) {
    std::vector<uint8_t> padded(block_size);
    padded[0] = 0x00;
    padded[1] = 0x02;
    
    // Fill with random non-zero bytes
    std::random_device rd;
    std::mt19937 gen(rd());
    std::uniform_int_distribution<> dis(1, 255);
    
    for (size_t i = 2; i < block_size - message.size() - 1; i++) {
        padded[i] = dis(gen);
    }
    
    padded[block_size - message.size() - 1] = 0x00;
    std::copy(message.begin(), message.end(), 
              padded.begin() + block_size - message.size());
    
    return padded;
}
\end{lstlisting}

\section{Performance Considerations}

\subsection{Computational Complexity}

\begin{itemize}
    \item \textbf{Key Generation:} $O(k^4)$ where $k$ is the key size in bits
    \item \textbf{Encryption/Decryption:} $O(k^3)$ for each operation
    \item \textbf{CRT Decryption:} Approximately 4x faster than standard decryption
\end{itemize}

\subsection{Optimization Techniques}

\begin{enumerate}
    \item \textbf{Use CRT for Decryption:} Can speed up decryption by 3-4x
    \item \textbf{Choose Optimal Public Exponent:} $e = 65537$ is commonly used
    \item \textbf{Use Efficient Big Integer Libraries:} GMP, OpenSSL, or Crypto++
    \item \textbf{Implement Montgomery Multiplication:} For very large numbers
\end{enumerate}

\section{Real-World Applications}

\subsection{Digital Signatures}

RSA can be used to create digital signatures:

\begin{lstlisting}[language=C++, caption=RSA Digital Signature]
class RSASignature {
private:
    RSA rsa;
    
public:
    // Sign a message
    uint64_t sign(uint64_t message_hash) {
        return rsa.decrypt(message_hash); // Use private key
    }
    
    // Verify a signature
    bool verify(uint64_t message_hash, uint64_t signature) {
        uint64_t decrypted = rsa.encrypt(signature); // Use public key
        return decrypted == message_hash;
    }
};
\end{lstlisting}

\subsection{Hybrid Encryption}

RSA is often used in hybrid systems:

\begin{enumerate}
    \item Generate a random symmetric key
    \item Encrypt the message with the symmetric key (AES, ChaCha20, etc.)
    \item Encrypt the symmetric key with RSA
    \item Send both the encrypted message and encrypted key
\end{enumerate}

\section{Conclusion}

RSA cryptography remains one of the most important and widely used asymmetric encryption algorithms. Understanding its mathematical foundations is crucial for implementing it correctly and securely.

\begin{prosbox}
\textbf{Advantages of RSA:}
\begin{itemize}
    \item Well-understood and extensively analyzed
    \item Supports both encryption and digital signatures
    \item Widely supported in cryptographic libraries
    \item Relatively simple to implement (basic version)
\end{itemize}
\end{prosbox}

\begin{consbox}
\textbf{Limitations of RSA:}
\begin{itemize}
    \item Slower than symmetric encryption
    \item Requires large key sizes for security
    \item Vulnerable to quantum attacks (Shor's algorithm)
    \item Complex to implement securely (padding, timing attacks)
\end{itemize}
\end{consbox}

\subsection{Future Considerations}

\begin{warningbox}
\textbf{Quantum Threat:} Shor's algorithm can factor large numbers efficiently on quantum computers, potentially breaking RSA. Post-quantum cryptography research is ongoing.
\end{warningbox}

\section{Reference Sources and Standards}

\subsection{Primary Research Papers}

\begin{enumerate}
    \item \textbf{Rivest, R. L., Shamir, A., \& Adleman, L.} (1978). "A Method for Obtaining Digital Signatures and Public-Key Cryptosystems." \textit{Communications of the ACM}, 21(2), 120-126. DOI: 10.1145/359340.359342
    
    \item \textbf{Rivest, R. L., \& Shamir, A.} (1982). "Efficient factoring based on partial information." \textit{Advances in Cryptology - EUROCRYPT '85}, 31-34.
    
    \item \textbf{Boneh, D.} (1999). "Twenty years of attacks on the RSA cryptosystem." \textit{Notices of the American Mathematical Society}, 46(2), 203-213.
\end{enumerate}

\subsection{Standards and Specifications}

\begin{enumerate}
    \item \textbf{RFC 8017} (2016). "PKCS \#1: RSA Cryptography Specifications Version 2.2." IETF. Available: \url{https://tools.ietf.org/html/rfc8017}
    
    \item \textbf{FIPS 186-4} (2013). "Digital Signature Standard (DSS)." NIST. Available: \url{https://nvlpubs.nist.gov/nistpubs/FIPS/NIST.FIPS.186-4.pdf}
    
    \item \textbf{NIST SP 800-56B} (2017). "Recommendation for Pair-Wise Key Establishment Schemes Using Integer Factorization Cryptography." NIST. Available: \url{https://nvlpubs.nist.gov/nistpubs/SpecialPublications/NIST.SP.800-56Br2.pdf}
    
    \item \textbf{NIST SP 800-57} (2020). "Recommendation for Key Management, Part 1: General." NIST. Available: \url{https://nvlpubs.nist.gov/nistpubs/SpecialPublications/NIST.SP.800-57pt1r5.pdf}
    
    \item \textbf{ANSI X9.31} (1998). "Digital Signatures Using Reversible Public Key Cryptography for the Financial Services Industry (rDSA)."
    
    \item \textbf{IEEE 1363} (2000). "Standard Specifications for Public-Key Cryptography."
\end{enumerate}

\subsection{Textbooks and Educational Resources}

\begin{enumerate}
    \item \textbf{Menezes, A. J., van Oorschot, P. C., \& Vanstone, S. A.} (1996). "Handbook of Applied Cryptography." CRC Press. ISBN: 0-8493-8523-7
    
    \item \textbf{Stinson, D. R.} (2005). "Cryptography: Theory and Practice." CRC Press. ISBN: 1-58488-508-4
    
    \item \textbf{Katz, J., \& Lindell, Y.} (2014). "Introduction to Modern Cryptography." CRC Press. ISBN: 1-4665-7026-1
    
    \item \textbf{Boneh, D., \& Shoup, V.} (2020). "A Graduate Course in Applied Cryptography." Available: \url{https://toc.cryptobook.us/}
    
    \item \textbf{Paar, C., \& Pelzl, J.} (2009). "Understanding Cryptography: A Textbook for Students and Practitioners." Springer. ISBN: 3-642-04100-0
\end{enumerate}

\subsection{Implementation Guides and Best Practices}

\begin{enumerate}
    \item \textbf{OpenSSL Documentation.} Available: \url{https://www.openssl.org/docs/}
    
    \item \textbf{Cryptographic Standards and Guidelines.} NIST Computer Security Resource Center. Available: \url{https://csrc.nist.gov/projects/cryptographic-standards-and-guidelines}
    
    \item \textbf{OWASP Cryptographic Storage Cheat Sheet.} Available: \url{https://cheatsheetseries.owasp.org/cheatsheets/Cryptographic_Storage_Cheat_Sheet.html}
    
    \item \textbf{Microsoft Security Development Lifecycle (SDL) Cryptographic Standards.} Available: \url{https://www.microsoft.com/en-us/securityengineering/sdl/cryptographic-standards}
\end{enumerate}

\section{Test Vectors and Validation}

\subsection{Standard Test Vectors}

\subsubsection{NIST Test Vectors}

NIST provides comprehensive test vectors for RSA implementations. These can be found in:

\begin{enumerate}
    \item \textbf{FIPS 186-4 Test Vectors:} \url{https://csrc.nist.gov/projects/cryptographic-algorithm-validation-program/digital-signatures}
    
    \item \textbf{NIST Cryptographic Algorithm Validation Program (CAVP):} \url{https://csrc.nist.gov/projects/cryptographic-algorithm-validation-program}
\end{enumerate}

\subsubsection{RFC 8017 Test Vectors}

RFC 8017 includes several test vectors for RSA operations:

\begin{examplebox}
\textbf{Example Test Vector from RFC 8017:}

\textbf{Key Generation:}
\begin{itemize}
    \item $p = 61$ (prime)
    \item $q = 53$ (prime)
    \item $n = p \times q = 3233$
    \item $\phi(n) = (p-1)(q-1) = 3120$
    \item $e = 17$ (public exponent)
    \item $d = 2753$ (private exponent)
\end{itemize}

\textbf{Encryption Test:}
\begin{itemize}
    \item Message: $M = 123$
    \item Ciphertext: $C = 123^{17} \bmod 3233 = 855$
\end{itemize}

\textbf{Decryption Test:}
\begin{itemize}
    \item Ciphertext: $C = 855$
    \item Decrypted: $M = 855^{2753} \bmod 3233 = 123$
\end{itemize}
\end{examplebox}

\subsection{OpenSSL Test Vectors}

OpenSSL provides test vectors in its test suite. You can find them in the OpenSSL source code:

\begin{lstlisting}[language=bash, caption=OpenSSL Test Vector Location]
# Clone OpenSSL repository
git clone https://github.com/openssl/openssl.git
cd openssl

# Test vectors are located in:
# test/rsa_test.c
# test/rsa_mp_test.c
# test/rsapss_test.c
\end{lstlisting}

\subsection{Where to Find Test Vectors}

\subsubsection{Online Resources}

\begin{enumerate}
    \item \textbf{NIST CAVP Database:} \url{https://csrc.nist.gov/projects/cryptographic-algorithm-validation-program/validation-search}
    
    \item \textbf{OpenSSL Test Suite:} \url{https://github.com/openssl/openssl/tree/master/test}
    
    \item \textbf{WolfSSL Test Vectors:} \url{https://github.com/wolfSSL/wolfssl/tree/master/certs}
    
    \item \textbf{Bouncy Castle Test Vectors:} \url{https://github.com/bcgit/bc-java/tree/master/core/src/test/java/org/bouncycastle/crypto/test}
    
    \item \textbf{Crypto++ Test Vectors:} \url{https://github.com/weidai11/cryptopp/tree/master/TestVectors}
\end{enumerate}

\subsubsection{Academic Sources}

\begin{enumerate}
    \item \textbf{University of California, Berkeley - Cryptography Course Materials:} \url{https://inst.eecs.berkeley.edu/~cs161/sp20/}
    
    \item \textbf{Stanford University - Applied Cryptography Group:} \url{https://crypto.stanford.edu/}
    
    \item \textbf{MIT OpenCourseWare - Cryptography and Cryptanalysis:} \url{https://ocw.mit.edu/courses/electrical-engineering-and-computer-science/6-875-cryptography-and-cryptanalysis-spring-2005/}
\end{enumerate}

\subsection{Validation Tools and Frameworks}

\subsubsection{Automated Testing}

\begin{enumerate}
    \item \textbf{Cryptographic Algorithm Validation System (CAVS):} NIST's official testing framework
    
    \item \textbf{OpenSSL Test Suite:} Comprehensive test suite for cryptographic implementations
    
    \item \textbf{WolfSSL Test Framework:} Automated testing for embedded systems
    
    \item \textbf{Bouncy Castle Test Framework:} Java-based cryptographic testing
\end{enumerate}

\subsubsection{Manual Validation}

\begin{lstlisting}[language=C++, caption=Manual Test Vector Validation]
#include <iostream>
#include <cassert>
#include "RSA.hpp"

void validate_test_vectors() {
    // Test vector from RFC 8017
    uint64_t p = 61;
    uint64_t q = 53;
    uint64_t n = 3233;
    uint64_t e = 17;
    uint64_t d = 2753;
    
    // Test encryption
    uint64_t message = 123;
    uint64_t expected_ciphertext = 855;
    
    uint64_t ciphertext = mod_pow(message, e, n);
    assert(ciphertext == expected_ciphertext);
    std::cout << "Encryption test passed!\n";
    
    // Test decryption
    uint64_t decrypted = mod_pow(ciphertext, d, n);
    assert(decrypted == message);
    std::cout << "Decryption test passed!\n";
}

int main() {
    validate_test_vectors();
    return 0;
}
\end{lstlisting}

\subsection{Performance Benchmarking}

\subsubsection{Benchmarking Tools}

\begin{enumerate}
    \item \textbf{OpenSSL Speed Test:} \texttt{openssl speed rsa}
    
    \item \textbf{Crypto++ Benchmarks:} Built-in benchmarking suite
    
    \item \textbf{WolfSSL Benchmarks:} Performance testing utilities
    
    \item \textbf{Custom Benchmarking Scripts:} Using high-resolution timers
\end{enumerate}

\subsubsection{Expected Performance}

\begin{table}[h]
\centering
\begin{tabular}{|c|c|c|c|}
\hline
\textbf{Key Size} & \textbf{Key Generation} & \textbf{Encryption} & \textbf{Decryption} \\
\hline
1024 bits & ~50ms & ~0.1ms & ~1ms \\
2048 bits & ~200ms & ~0.3ms & ~5ms \\
4096 bits & ~800ms & ~1ms & ~20ms \\
\hline
\end{tabular}
\caption{Typical RSA Performance (Intel i7, single-threaded)}
\end{table}

\subsection{Further Reading}

\begin{itemize}
    \item "Handbook of Applied Cryptography" by Menezes, van Oorschot, and Vanstone
    \item "Cryptography: Theory and Practice" by Stinson
    \item RFC 8017: PKCS\#1 v2.2
    \item NIST Special Publication 800-56B: Key Establishment Schemes
    \item "A Graduate Course in Applied Cryptography" by Boneh and Shoup
    \item "Understanding Cryptography" by Paar and Pelzl
\end{itemize}

\appendix

\section{Mathematical Proofs}

\subsection{Proof of Euler's Theorem}

\begin{proof}
Let $S = \{a_1, a_2, \ldots, a_{\phi(n)}\}$ be the set of integers coprime to $n$.

Consider the set $T = \{aa_1, aa_2, \ldots, aa_{\phi(n)}\}$ where $\gcd(a,n) = 1$.

We can show that:
\begin{enumerate}
    \item All elements in $T$ are coprime to $n$
    \item All elements in $T$ are distinct modulo $n$
    \item Therefore, $T$ is a permutation of $S$ modulo $n$
\end{enumerate}

This implies:
$$(aa_1)(aa_2)\cdots(aa_{\phi(n)}) \equiv a_1a_2\cdots a_{\phi(n)} \pmod{n}$$

Since $\gcd(a_i, n) = 1$ for all $i$, we can cancel the $a_i$ terms:
$$a^{\phi(n)} \equiv 1 \pmod{n}$$
\end{proof}

\section{Code Examples}

\subsection{Complete RSA Implementation}

The complete implementation is available in the CryptoGL library at \texttt{src/RSA.cpp} and \texttt{src/RSA.hpp}.

\subsection{Testing RSA Implementation}

\begin{lstlisting}[language=C++, caption=RSA Test Suite]
#include <cassert>
#include <iostream>

void test_rsa() {
    // Test with small numbers
    RSA rsa_small(64);
    
    uint64_t test_message = 12345;
    uint64_t encrypted = rsa_small.encrypt(test_message);
    uint64_t decrypted = rsa_small.decrypt(encrypted);
    
    assert(test_message == decrypted);
    std::cout << "Small RSA test passed!\n";
    
    // Test with larger numbers
    RSA rsa_large(512);
    
    uint64_t large_message = 987654321;
    uint64_t large_encrypted = rsa_large.encrypt(large_message);
    uint64_t large_decrypted = rsa_large.decrypt(large_encrypted);
    
    assert(large_message == large_decrypted);
    std::cout << "Large RSA test passed!\n";
}

int main() {
    test_rsa();
    return 0;
}
\end{lstlisting}

\end{document} 